\documentclass[11pt]{article}

\usepackage[utf8]{inputenc}
\usepackage[T1]{fontenc}
\usepackage{amsmath, amssymb}
\usepackage{geometry}
\usepackage{hyperref}

\geometry{margin=1in}

\title{Collected Abstracts}
\newcommand{\authorname}{Reivax Del Ponte}
\newcommand{\institute}{Pseudo-Institute of Non-Experimental Physics}
\author{\authorname \thanks{\institute}}

\date{\today}

\begin{document}
\maketitle

\section*{Article 1: Photons and the Geometry of Spacetime}
\subsection*{Reasoning About Photon Reference Frames}

\begin{abstract}
We examine the consequences of taking the reference frame of a single photon seriously, treating emission and absorption as coincident events in that degenerate limit, and explore how this perspective informs our reasoning about multiple photons, the structure of spacetime, and the topology of causal sequences. We argue that, while the photon's proper frame collapses time and distance, the topology of spacetime is preserved, and a consistent picture of how photons and event sequences are to be reasoned about within that structure can be reconstructed.
\end{abstract}

\section*{Article 2: A Photon-Frame Theory}
\subsection*{Spacetime as a Lightlike Graph}

\begin{abstract}
We propose to take seriously the perspective that, in a physical theory, the only legitimate reference frames are those of photons themselves. The standard objection to this position --- that Lorentz boosts to $v = c$ are singular --- is set aside as a feature rather than a bug: the singularity encodes the fact that emission and absorption become coincident in the photon's frame. Under this restriction, ordinary spacetime coordinates are no longer available, and the only geometric structure left is a graph whose edges are null geodesics (photon worldlines) and whose vertices are emission and absorption events. We call this structure the \emph{lightlike graph}. We show that, when three long-standing problems --- Bell's theorem and quantum nonlocality, the measurement problem of quantum mechanics, and the black-hole information paradox --- are recast on this graph, several apparent contradictions soften into questions about graph topology, and the resulting reformulations suggest fresh avenues of attack.
\end{abstract}

\section*{Article 3: Toward a Unified Theory}
\subsection*{Schr\"odinger, Spacetime, and the Lightlike Graph}

\begin{abstract}
Quantum mechanics and general relativity rest on incompatible pictures of time. The time-dependent and time-independent Schr\"odinger equations treat time as a distinguished external parameter, while general relativity binds time to geometry in the form of spacetime. The mismatch is usually papered over by treating quantum mechanics as a theory of states on a fixed spacetime background, but the cost of doing so is the open problem of wave-function collapse and the unresolved status of the ``observer''. We propose that, as in the companion papers, restricting attention to the lightlike graph softens the mismatch: the two Schr\"odinger equations both become statements about graph incidence --- one about the existence of an edge, the other about the topology of its endpoints --- and the observer becomes a vertex in the graph rather than a frame-bound external agent. We argue that this is the right vantage point from which to begin searching for a unification.
\end{abstract}

\section*{Article 4: A Lightlike-Graph Audit of String Theory}
\subsection*{Pseudo-Science, Productive Mathematics, or Something Else?}

\begin{abstract}
A vigorous public dispute surrounds string theory. Critics (most prominently Smolin and Woit) argue that, after four decades of effort, it has produced no experimental prediction sharp enough to falsify it and has therefore forfeited the status of a scientific theory. Defenders argue that the theory's mathematical coherence, its many dualities, and its unification of gravity with the other interactions constitute progress that should not be judged by the standards of an earlier, more empirically forgiving, era. We do not take sides. Instead, we audit the dispute using the lightlike graph $\mathcal{G}$ introduced in the companion papers, in which the only primitive is the lightlike relation between emission and absorption events and the only questions are combinatorial: which edges exist, which vertices they share, how the degrees organise themselves. The audit is a diagnostic, not a verdict: it asks whether the apparatus of string theory produces new combinatorial possibilities on $\mathcal{G}$ (and would therefore generate new knowledge), exposes combinatorial redundancies already accessible (and would therefore be unnecessary complication), or misses combinatorial structure that the more austere setting makes visible (and would therefore suggest a redirection of effort). Each of the three outcomes is logically possible, and a preliminary inspection suggests that string theory's status on $\mathcal{G}$ depends on whether one restricts attention to kinematic structure or to dynamical content.
\end{abstract}

\section*{Article 5: Source-Tagged Lightlike Graphs}
\subsection*{When the Extra Bookkeeping Helps, and When It Doesn't}

\begin{abstract}
In the companion paper auditing string theory on the lightlike graph $\mathcal{G}$, we noted that the combinatorial content of string-theoretic dualities (and, conditionally, of compactifications) is visible on $\mathcal{G}$ only under one caveat: \emph{if one knows which photons were emitted by which source}. That caveat identified a piece of additional bookkeeping that lies outside the bare graph and that may, in several places, carry genuine combinatorial work. In this paper we give that bookkeeping a name, formalise it as a labelling of the lightlike graph, and audit its productivity in the same spirit as the previous paper. The bookkeeping is a labelling of edges by the source that emitted them (and, optionally, by the entangled pair to which they belong, by the object on which they were absorbed, and by the spacetime region in which they originated). We examine this labelling on the problems previously cast on $\mathcal{G}$ --- Bell's theorem, the measurement problem, the two-photon meeting, the black-hole information paradox, the sequence problem, and the T-duality situation --- and in each case ask whether the added labels do visible combinatorial work. The conclusion is mixed: the bookkeeping is genuinely productive when many photons share an emission or absorption vertex and disambiguation matters; it is unproductive (and possibly harmful) when the source is already identifiable from topology alone or when the photons are physically indistinguishable. The bookkeeping is best understood not as a universal extension of $\mathcal{G}$ but as an opt-in refinement, used where the bare graph is ambiguous.
\end{abstract}

\section*{Article 6: The Quantum Eraser on the Lightlike Graph}
\subsection*{Does Photon-Frame Combinatorics Resolve a Paradox, or Merely Rephrase It?}

\begin{abstract}
The delayed-choice quantum eraser exhibits a striking asymmetry of temporal ordering: a measurement choice made on an idler photon strictly \emph{after} its entangled signal partner has been recorded appears to determine whether the signal history belongs to an interference pattern or to a particle-like distribution. We project the experiment onto the lightlike graph $\mathcal{G}$ developed in the companion papers, in which emission and absorption coincide as vertices and photon worldlines are directed edges. Two bookkeeping augmentations from article~5 --- pair tagging and a new \emph{future-tagging} of branched absorption vertices --- turn out to do genuine combinatorial work, and we argue that the apparent paradox softens into a question about which of several discrete graph completions is recorded as ``the'' completion. The resolution is partial: lightlike graph theory rephrases the question, forecloses the retrocausality reading, and brings the bookkeeping that the experiment actually uses (entangled pairs, delayed coincidence) into the formal vocabulary, but it leaves untouched the deeper question of why one completion is selected rather than another. The lightlike graph is therefore not a resolution of the paradox but a clarification of what the paradox was asking.
\end{abstract}

\section*{Article 7: Time Crystals on the Lightlike Graph}
\subsection*{Periodicity Without Proper Time}

\begin{abstract}
A time crystal is, in the original Wilczek sense, a system whose lowest-energy state breaks time-translation symmetry: the ground state has a periodic structure that no translation in time can erase. Subsequent work has distinguished this proposal from the now experimentally realised \emph{discrete time crystals}, in which a periodically driven many-body system responds at a subharmonic frequency of the drive. We project the phenomenon onto the lightlike graph $\mathcal{G}$ developed in the companion papers. The restriction to photon frames has already discarded proper time, so the notion of ``oscillation in time'' has no direct representation; what we have instead is a \emph{recurrence} of certain subgraphs across a sequence of lightlike graphs, and an \emph{identifiability} question for repeated patterns. We find that the lightlike graph captures the empirical fingerprint of time-crystalline behaviour (subharmonic response, rigidity under perturbations, the role of a drive photon) in purely combinatorial terms, and that the bookkeeping of articles~5 and~6 --- source tagging, region tagging, future tagging --- can be applied with pay-off. The clarification we obtain is partial: the graph records what recurrence is, forecloses several naive readings (the system is not ``moving'' in any proper time), and provides a natural home for the distinction between equilibrium time crystals (forbidden by the standard no-go theorems) and driven (Floquet) time crystals (allowed and observed). The deeper questions --- why subharmonic responses are stable, and what selects the period --- remain, as in the previous papers, outside the combinatorics.
\end{abstract}

\section*{Article 8: Scalar Quantum Field Theory on the Lightlike Graph}
\subsection*{What the Simplest Possible Field Looks Like When Spacetime Has Been Discarded}

\begin{abstract}
Quantum field theory, in its most elementary form, is the theory of a single real scalar field $\phi(x)$ --- a number assigned to every point $x$ of spacetime --- and a Lagrangian density $\mathcal{L}[\phi, \partial\phi]$ built from it. The picture is intuitive: a field is a height-map over a fixed landscape, and excitations are bumps that move on it. The intuition weakens in general relativity, where the landscape is curved and the height-map stands on a dynamic geometry whose own excitations are governed by a different field. We project the scalar-field picture onto the lightlike graph $\mathcal{G}$ developed in the companion papers, in which spacetime has been discarded and the only primitive is the lightlike relation between emission and absorption events. The result is, on first inspection, very unintuitive: there is no landscape to stand on, no point $x$ to evaluate $\phi$ at, and no derivative to feed into $\mathcal{L}$. We propose three strategies for dealing with this unintuitiveness --- re-reading the field as a labelling of the corpus of photon edges, re-reading excitations as differences between successive graphs in the article-3 sequence, and re-reading selection rules as admissible-graph constraints --- and find that, while the questions the conventional scalar QFT asks are not all answerable in the photon-frame setting, the \emph{clarifying} work that the lightlike graph does for simpler puzzles transfers directly. The scalar field on $\mathcal{G}$ is not a different theory; it is the same theory with a different vocabulary, and the untranslatable residue is, once again, the dynamical content (vacuum energy, renormalisation, symmetry breaking) rather than the bookkeeping of which field value sits where.
\end{abstract}

\end{document}
